\section{Establishing a trustworthy simulation}

\subsection{From a proof of concept to shared excitation physics}

The June implementation explored a decoupled illumination-design loss, avoiding differentiation through the full iterative reconstruction. It used a bimodal intensity target and axial/mirror decorrelation, initially with a scalar pupil-focusing model. Later objectives used exponential speckle statistics. Training reduced mirror correlation and showed that changing phase structure could break depth symmetry. This established a useful design direction, but not the final forward model.

The September physics audit found that the excitation paths did not share the same model: the imported DMD path used a vectorial PSF, whereas the phase-encoder path used scalar propagation, and the initial new coder used an objective-focusing geometry rather than the intended relay. FRIS replaced these with a common coherent vectorial relay. Intensity is formed after propagating the complex field and summing the squared magnitudes of its vector components; emission is modelled incoherently.

The audit also aligned the excitation focus with array index 15 on the 32-plane analysis grid. Mirror pairs must be defined about that focus rather than the midpoint of an even-length array. These corrections made encoder comparisons meaningful and superseded numerical claims from the earlier model.

\subsection{Phase-feature scale and relay geometry}

Early phase optimisation used 0.125 \ensuremath{\mu}m features. Re-evaluating the designed element on a 0.32 \ensuremath{\mu}m phase grid weakened its apparent mirror-separation advantage. The practical conclusion was to prevent training from relying on unresolved or non-propagating fine structure and to separate physical phase features from the finer numerical propagation grid. The 0.32 \ensuremath{\mu}m feature grid and 0.08 \ensuremath{\mu}m optical grid became the later reference.

The output-hop audit identified a 3.3 nm value in the inherited configuration, inconsistent with micrometre layer spacings. A 6.6 \ensuremath{\mu}m final gap was adopted after a gap study balanced angular-pattern changes against light leaving the finite window. The logged random two-layer study found near-angle correlation decreasing from 0.62 to 0.34 and 0.29 at gaps of 0, 6.6 and 10 \ensuremath{\mu}m, while relative retained energy fell from 1.00 to 0.73 and 0.63. These are historical two-layer diagnostics, not final one-plate recovery scores.

The implementation was later expressed as a phase element near an intermediate image plane. Fresnel scaling with relay demagnification 60 gives approximately 19 \ensuremath{\mu}m phase features, 24 mm defocus and 8--14 mrad element-side tilts. Numerical paraxial controls supported this interpretation; no hardware implementation was tested.

\subsection{Fine measurements, integrated voxels and fast construction}

Measurements are generated from a fine object; the reconstruction reference is its voxel average. They are not generated by simply replicating the coarse reference. The integrated operator is \(\mathsf A=\mathsf B\mathsf U\), where \(\mathsf B\) is the fine measurement process and \(\mathsf U\) replicates each reconstruction value over its fine patch.

The FYP derivation was restored in the October implementation: a fine emitter's response is its illumination value multiplied by a reusable shifted emission footprint. Summing these products over a voxel patch constructs the integrated operator without repeatedly performing volumetric convolutions. The final one-plate build was logged at 109 s including preparation and validation. Live-forward checks were of order \(8\times10^{-7}\) relative error. A historical two-layer brute build took about 15 hours, but the configurations and timing boundaries were not matched, so their ratio is not a controlled speed-up measurement.

The retained patch-versus-coarse mismatch test showed 1.07\% versus 25.7\% forward discrepancy and 0.177 versus 1.859 \ensuremath{\mu}m depth error on subvolume 0. Direct optical simulation at 0.64 \ensuremath{\mu}m undersamples the high-NA field. This established that optical detail should be retained before integration onto the reconstruction basis. The auxiliary point-sampled arm is outside the accepted main comparison and is not needed for the conclusions here.
